\chapter{Pilonidal Sinus Surgery}

\section*{Introduction \& Clinical Indications}
Pilonidal sinus surgery is a common surgical intervention performed to treat pilonidal disease, a chronic inflammatory condition characterized by the formation of a sinus tract or abscess in the natal cleft, the groove between the buttocks. The term "pilonidal," derived from Latin "pilus" (hair) and "nidus" (nest), accurately reflects the etiology, as embedded hairs and skin debris are central to its pathogenesis. This procedure aims to definitively remove the diseased tissue, including the sinus tract, any associated pits, and foreign material, thereby alleviating symptoms such as pain, swelling, and recurrent discharge of pus or blood. It is a crucial treatment for patients experiencing recurrent infections, chronic drainage, or significant discomfort that impacts their quality of life, offering a permanent solution where conservative measures have failed.

\begin{itemize}
  \item Pilonidal disease surgery
  \item Pilonidal cystectomy
  \item Pilonidal sinus excision
  \item Cleft lift procedure
  \item Pilonidal abscess incision and drainage with later excision
  \item Karydakis flap procedure
\end{itemize}

The fundamental surgical principle for pilonidal disease is the complete eradication of the pathological sinus tract, its associated pits, and all embedded foreign material, primarily hair and keratinous debris. This foreign-body reaction and chronic inflammation create an epithelialized or granulation tissue-lined tract that acts as a nidus for recurrent infection. By removing this entire diseased complex, the surgeon eliminates the source of chronic inflammation and potential for bacterial colonization, thereby promoting definitive healing.

The rationale behind various surgical approaches centers on not only excising the disease but also managing the resulting wound to minimize recurrence and optimize healing. Techniques range from simple excision with open healing (secondary intention), which prioritizes low recurrence at the cost of prolonged wound care, to primary closure, which offers faster healing but a higher recurrence risk. More advanced flap procedures, such as the Karydakis or Bascom cleft lift, aim to flatten the natal cleft and shift the incision line away from the midline, thereby reducing shear forces, moisture accumulation, and the likelihood of hair ingress, which are key factors in disease recurrence. The choice of technique is tailored to the extent of disease, presence of acute infection, and patient factors.

Pilonidal disease, though likely present for centuries, gained significant medical attention in the 19th century. Herbert Mayo, in 1833, is often credited with one of the earliest descriptions. However, its prevalence surged during World War II, particularly among military personnel who spent long hours driving jeeps, leading to the colloquial term "jeep driver's disease." This era highlighted the role of prolonged sitting, friction, and trauma in its pathogenesis.

Early surgical approaches, predominantly wide excision with healing by secondary intention, were championed by surgeons like Patey and Scarff in the mid-20th century. While effective in reducing recurrence, these methods often resulted in large, slow-healing wounds requiring extensive postoperative care. The late 20th century saw a paradigm shift towards techniques designed to close the wound and modify the natal cleft anatomy. The Karydakis flap, introduced by Karydakis in 1973, and the Bascom cleft lift procedure, popularized by Bascom in the 1980s, revolutionized pilonidal surgery by demonstrating significantly lower recurrence rates and faster healing through off-midline closures and cleft flattening. These advancements continue to influence modern surgical practice, emphasizing anatomical modification to prevent recurrence.

Surgical intervention for pilonidal disease is indicated in patients presenting with symptomatic or recurrent disease that significantly impacts their quality of life. The primary clinical indications include:

\begin{itemize}
  \item Recurrent pilonidal sinus disease with repeated episodes of infection, pain, and purulent or serosanguinous drainage.
  \item A chronically draining sinus tract that fails to heal despite conservative measures such as meticulous hygiene, hair removal, and antibiotic therapy for acute flares.
  \item Persistent pain, swelling, or discomfort in the sacrococcygeal region caused by the pilonidal sinus, especially when exacerbated by sitting or physical activity.
  \item An acutely infected pilonidal abscess requiring incision and drainage, with subsequent planned definitive excision to prevent recurrence.
  \item Extensive or complex pilonidal disease, including multiple sinus openings or tracts, that is unlikely to resolve with non-operative management.
  \item Disease that interferes with daily activities, work, or personal hygiene, leading to significant functional impairment.
  \item Cosmetic concerns or psychological distress related to the chronic drainage, odor, or appearance of the pilonidal lesion.
\end{itemize}

Pilonidal sinus surgery holds a dual clinical positioning within the management algorithm for pilonidal disease. It serves as the primary definitive treatment for established, symptomatic, and recurrent pilonidal disease, particularly when conservative measures have proven ineffective. In these scenarios, surgery aims for complete cure and prevention of future episodes.

Conversely, surgery can also be considered a secondary or staged procedure. For instance, an acute pilonidal abscess is initially managed with incision and drainage to relieve pain and control infection; definitive excisional surgery is then typically deferred until the acute inflammation has subsided. Furthermore, for patients with incidental, asymptomatic pilonidal pits, or those with very mild, first-time presentations, conservative management focusing on hygiene, hair removal, and lifestyle modifications is often the first-line approach. Surgery is thus reserved for those who fail conservative treatment, experience recurrent symptoms, or present with complicated disease.

\section*{Relevant Surgical Anatomy}
Pilonidal disease primarily affects the sacrococcygeal region, specifically the natal cleft, which is the deep furrow separating the buttocks. This area is characterized by several anatomical features that predispose to the condition. The skin in the natal cleft is often hirsute, with coarse hairs that can become ingrown or penetrate the skin due to friction and pressure, particularly during prolonged sitting. Beneath the skin lies subcutaneous fat, which can become extensively involved in chronic pilonidal disease, forming complex sinus tracts and abscess cavities.

Deep to the subcutaneous tissue, the sacrococcygeal fascia covers the sacrum and coccyx. The gluteus maximus muscles originate from the sacrum and coccyx and insert into the femur, forming the bulk of the buttocks and contributing to the depth of the natal cleft. The coccyx itself is the terminal segment of the vertebral column, and while not directly involved in the disease process, its proximity defines the anatomical boundaries of the surgical field. Understanding the depth and lateral extent of the sinus tracts relative to these structures is critical for complete excision and successful reconstruction, especially in complex or recurrent cases.

\section*{Preoperative Workup \& Preparation}
Thorough preoperative workup and preparation are essential to optimize outcomes and minimize complications in pilonidal sinus surgery. This typically begins with a detailed medical history and physical examination to assess the patient's overall health status and the extent of the pilonidal disease. Any acute infection or abscess must be addressed first; an actively inflamed, fluctuant abscess should be incised and drained, with definitive excision deferred until the acute process has resolved.

Standard preoperative laboratory tests, such as a complete blood count and basic metabolic panel, are usually performed. Imaging studies are generally not required for routine pilonidal disease but may be considered for very complex or recurrent cases to delineate the extent of the tracts or rule out other pathologies. Patients are instructed to be NPO (nil per os) for at least 6-8 hours prior to surgery in preparation for general anesthesia. Medications, particularly anticoagulants or antiplatelet agents, are reviewed and adjusted as per institutional guidelines to minimize bleeding risk.

Prophylactic broad-spectrum antibiotics, typically a single dose administered intravenously prior to incision, are often given to reduce the risk of surgical site infection, especially for closed wounds. The surgical site may be shaved or clipped in the operating room to remove hair from the natal cleft. Comprehensive informed consent is obtained, detailing the chosen surgical technique, potential risks, benefits, and the expected postoperative course, including the implications of different wound management strategies (e.g., open vs. closed healing, recurrence rates). Patients are also advised on meticulous hygiene and hair control measures to be continued postoperatively.

While pilonidal sinus surgery is generally safe and effective, certain conditions may contraindicate or necessitate precautions for the procedure.

\textbf{Situations requiring conservative care, infection treatment, optimization, or an alternative plan:}
\begin{itemize}
  \item An acute, undrained pilonidal abscess with active cellulitis or systemic signs of infection is a contraindication to immediate definitive excision with primary closure. It requires initial incision and drainage to control the infection, with elective surgery deferred.
  \item Uncontrolled systemic infection or sepsis, which must be managed and stabilized prior to any elective surgical intervention.
  \item Severe, uncorrectable coagulopathy that poses an unacceptable risk of intraoperative or postoperative hemorrhage.
\end{itemize}

\textbf{Relative Contraindications \& Precautions:}
\begin{itemize}
  \item Significant comorbidities (e.g., severe cardiac, pulmonary, or renal disease) that substantially increase the risks associated with general or regional anesthesia. These patients require thorough preoperative optimization.
  \item A small, asymptomatic pilonidal pit without a history of infection or drainage, which can often be managed conservatively with hygiene and hair removal.
  \item Immunocompromised states (e.g., diabetes, HIV, immunosuppressive therapy) may increase the risk of wound infection and delayed healing, requiring careful perioperative management and antibiotic prophylaxis.
  \item Previous failed pilonidal surgery, which may necessitate a more complex reconstructive approach and careful planning.
\end{itemize}
Precautions also include meticulous attention to hair control in the natal cleft, both pre- and postoperatively, and careful patient counseling regarding wound care and activity restrictions to optimize healing and prevent recurrence.

\section*{Surgical Technique \& Procedure}
Pilonidal sinus surgery is typically performed under general anesthesia, though regional anesthesia (e.g., spinal or epidural) can be an option for select patients. The patient is positioned prone (face down) on the operating table, often with the buttocks taped apart to maximize exposure of the natal cleft. The sacrococcygeal area is then meticulously prepared with an antiseptic solution and draped in a sterile fashion.

The specific surgical technique varies based on the extent of the disease and the surgeon's preferred method of wound management. Common approaches include simple excision with open healing, primary closure, or various flap procedures. A general procedural outline is as follows:

\begin{itemize}
  \item \textbf{Identification and Marking:} The surgeon identifies all visible sinus pits and tracts, often injecting methylene blue dye or similar agent into the pits to delineate the full extent of the disease process beneath the skin.
  \item \textbf{Excision:} An elliptical incision is made encompassing all identified pits and the underlying sinus tract. The diseased tissue, including skin, subcutaneous fat, and the entire sinus complex, is sharply dissected down to the presacral fascia, ensuring complete removal of all hair and granulation tissue.
  \item \textbf{Hemostasis:} Meticulous hemostasis is achieved using electrocautery to minimize postoperative bleeding and hematoma formation.
  \item \textbf{Wound Management (Choice of Closure):}
  \begin{itemize}
    \item \textbf{Open Excision (Healing by Secondary Intention):} The wound is left open to heal from the base upwards. This method has a lower recurrence rate but requires prolonged wound care.
    \item \textbf{Primary Closure:} The wound edges are approximated and sutured in layers. This offers faster healing but may have a higher recurrence rate, especially if the natal cleft is not flattened.
    \item \textbf{Flap Procedures (e.g., Karydakis, Cleft Lift):} These involve excising the sinus and then mobilizing a skin flap (often off-midline) to close the defect and flatten the natal cleft. This technique aims to reduce recurrence by altering the anatomy and shifting the incision away from the midline.
  \end{itemize}
  \item \textbf{Drain Placement:} For larger excisions or flap procedures, a suction drain may be placed in the wound cavity to prevent seroma or hematoma formation, typically removed within 24-48 hours.
  \item \textbf{Closure:} If closed, the skin is approximated with sutures or staples. For flap procedures, the skin is meticulously closed to ensure proper flap viability and tension.
  \item \textbf{Dressing:} A sterile dressing is applied to the wound.
\end{itemize}

The duration of the operation typically ranges from 30 minutes to 90 minutes, depending on the complexity of the disease and the chosen surgical technique. The choice of wound closure significantly influences the postoperative recovery period and the long-term recurrence risk.

\section*{Complications \& Risks}
As with any surgical procedure, pilonidal sinus surgery carries potential risks and complications, which can be broadly categorized as general surgical risks and procedure-specific risks.

\textbf{General Surgical Complications:}
\begin{itemize}
  \item \textbf{Bleeding/Hematoma:} Intraoperative or postoperative hemorrhage, potentially requiring re-operation.
  \item \textbf{Infection:} Surgical site infection (SSI), ranging from superficial cellulitis to deep abscess formation, particularly in closed wounds.
  \item \textbf{Anesthesia Risks:} Adverse reactions to anesthetic agents, respiratory or cardiovascular complications, and postoperative nausea and vomiting.
  \item \textbf{Pain:} Postoperative pain, which is usually manageable with analgesics.
  \item \textbf{Scarring:} Formation of a visible or hypertrophic scar.
\end{itemize}

\textbf{Procedure-Specific Complications:}
\begin{itemize}
  \item \textbf{Recurrence:} The most significant and common complication, with rates varying widely depending on the surgical technique (e.g., higher with primary closure, lower with flap procedures). Recurrence can manifest as new pits, drainage, or abscess formation.
  \item \textbf{Wound Dehiscence/Breakdown:} Particularly in closed wounds, leading to delayed healing or requiring conversion to an open wound.
  \item \textbf{Delayed Wound Healing:} Especially common with open excision techniques, where healing by secondary intention can take several weeks to months.
  \item \textbf{Seroma Formation:} Accumulation of serous fluid in the wound cavity, which may require aspiration or drainage.
  \item \textbf{Chronic Pain:} Persistent discomfort or neuropathic pain in the surgical area, though rare.
  \item \textbf{Flap Necrosis:} A specific risk of flap procedures, where compromised blood supply to the flap can lead to tissue death and wound breakdown.
  \item \textbf{Altered Sensation:} Numbness or altered sensation around the surgical site due to nerve irritation or transection.
  \item \textbf{Cosmetic Deformity:} Unsatisfactory aesthetic outcome, particularly with wide excisions or poorly executed flap closures.
\end{itemize}

\section*{Postoperative Care \& Recovery}
Immediately following pilonidal sinus surgery, patients are transferred to the Post-Anesthesia Care Unit (PACU) for close monitoring of vital signs, pain levels, and any signs of bleeding. Once stable, they are discharged home, often on the same day for uncomplicated cases, or admitted for overnight observation, particularly after more extensive procedures or if drains are in place. During the first 24-48 hours, pain management is a priority, typically with oral analgesics. Patients are advised to minimize prolonged sitting and to keep the surgical area clean and dry. For open wounds, initial dressing changes are often performed by nursing staff or taught to the patient/caregiver.

Over the first 1-2 weeks, recovery largely depends on the wound closure method. Patients with primary closure or flap procedures generally experience faster initial healing, with sutures or staples typically removed around 10-14 days post-op. Activity restrictions, such as avoiding strenuous exercise or heavy lifting, are usually in place for 2-4 weeks. For open wounds, healing by secondary intention can take anywhere from 6 weeks to several months, requiring diligent daily wound care, including packing and dressing changes. Long-term recovery involves maintaining excellent hygiene, regular hair removal in the natal cleft (e.g., shaving, laser hair removal), and avoiding prolonged sitting to minimize the risk of recurrence. Full return to unrestricted activities and complete wound maturation can take several months.

During pilonidal sinus surgery, the surgeon meticulously documents the extent of the disease, including the number and location of sinus pits, the length and complexity of the tracts, the presence of any abscess cavities, and the amount of embedded hair and debris. The excised tissue, comprising skin, subcutaneous fat, and the entire sinus tract, is routinely sent for histopathological examination. The pathology report typically confirms the diagnosis of pilonidal disease, characterized by chronic inflammation, foreign-body giant cell reaction to hair shafts, and often an epithelialized sinus tract or granulation tissue. The report also serves to exclude other rare conditions that can mimic pilonidal disease, such as hidradenitis suppurativa, anal fistula, or malignancy, although the latter is exceedingly rare in this context.

A normal and successful postoperative course following pilonidal sinus surgery is characterized by progressive wound healing, resolution of pain, and cessation of drainage. Patients can expect initial discomfort, which gradually subsides over days to weeks with appropriate pain medication. For closed wounds, the incision should appear clean and dry, with minimal swelling or redness, and sutures/staples removed at the scheduled follow-up. For open wounds, the wound bed should show healthy granulation tissue filling the defect from the base upwards, with decreasing exudate. The ultimate goal is complete epithelialization of the wound, leading to a healed, symptom-free natal cleft. Patients should experience a gradual return to normal activities, with the ability to sit comfortably and maintain personal hygiene without difficulty or recurrence of symptoms.

Postoperative care is crucial for optimal healing and prevention of recurrence. Patients are typically scheduled for several follow-up visits with the surgeon to monitor wound healing and address any concerns. Key aspects of postoperative care include:

\begin{itemize}
  \item \textbf{Wound Care:} Regular dressing changes, which may involve packing for open wounds, until complete healing is achieved. Patients or caregivers are educated on proper wound care techniques.
  \item \textbf{Suture/Staple Removal:} Scheduled removal of sutures or staples, usually 10-14 days post-surgery for closed wounds.
  \item \textbf{Hygiene and Hair Control:} Meticulous hygiene of the natal cleft and ongoing hair removal (e.g., shaving, depilatory creams, or laser hair removal) are strongly advised to prevent new hair ingress and recurrence.
  \item \textbf{Activity Restrictions:} Gradual return to normal activities, with avoidance of prolonged sitting, strenuous exercise, and activities that put direct pressure on the surgical site for several weeks.
  \item \textbf{Warning Signs:} Patients are educated on warning signs of complications, such as increasing pain, redness, swelling, fever, or purulent discharge, and instructed to contact their surgeon immediately if these occur.
\end{itemize}
Long-term follow-up may be less frequent once the wound is fully healed, but continued adherence to hygiene and hair control measures is paramount for sustained success.

\section*{Outcomes \& Prognosis}
Pilonidal disease is a relatively common condition, primarily affecting young adults, with an estimated incidence of 26 per 100,000 population. There is a strong male predominance, with males affected 2-4 times more frequently than females. The peak incidence occurs in the second and third decades of life, typically between the ages of 15 and 30, with a decline in incidence after age 40. Risk factors include male gender, obesity, hirsutism (especially coarse body hair), a deep natal cleft, prolonged sitting, sedentary lifestyle, and a family history of the disease. While not associated with mortality, pilonidal disease carries significant morbidity, leading to chronic pain, recurrent infections, time off work or school, and a substantial impact on quality of life. Recurrence rates after surgery vary widely depending on the technique, contributing to the burden of repeated healthcare encounters.

The outcomes and prognosis following pilonidal sinus surgery are generally favorable, with most patients achieving complete healing and symptomatic relief, though recurrence remains a significant concern. Success rates, defined as complete healing without recurrence, vary considerably based on the surgical technique employed. Open excision (healing by secondary intention) typically boasts recurrence rates between 5-15\% but with prolonged healing times. Primary closure offers faster healing but often higher recurrence rates, sometimes up to 20-30\%. Flap-based procedures, such as the Karydakis flap or Bascom cleft lift, have demonstrated superior outcomes with recurrence rates as low as 1-5\% in experienced hands, coupled with faster healing and better patient satisfaction. Prognostic factors for successful outcomes include complete excision of all diseased tissue, effective wound management, meticulous postoperative hygiene, and diligent hair control in the natal cleft. While no single landmark trial dictates management like in some other surgical fields, numerous comparative studies consistently support the efficacy of cleft-lift procedures in reducing recurrence. Overall, with appropriate surgical technique and diligent postoperative care, the long-term prognosis for patients undergoing pilonidal sinus surgery is good, allowing for a return to normal functional activities.

\section*{References}
\begin{enumerate}
  \item MedlinePlus. Pilonidal sinus surgery. U.S. National Library of Medicine; 2025.
  \item Steele SR, Hull TL, Read TE, et al. The ASCRS Textbook of Colon and Rectal Surgery. Springer; 3rd edition; 2016.
  \item American Society of Colon and Rectal Surgeons. Clinical Practice Guidelines for the Management of Pilonidal Disease. Dis Colon Rectum; 2019.
  \item Current Medical Diagnosis and Treatment. McGraw-Hill; 2025.
\end{enumerate}

\clearpage
